\begin{ouverture}[son]{15}{Période et perception sonore}{Ondes et signaux}{Observer · Modéliser · Expérimenter · Expliquer}
\objectif{Déterminer une période et une fréquence.}
\objectif{Distinguer hauteur, timbre et niveau sonore.}
\objectif{Lire une échelle non linéaire et identifier des situations à risque.}
\prerequis{Proportionnalité, conversions et lecture de graphiques.}
\end{ouverture}
\essentiel
\essSection{1}{Un motif qui se répète}
Un signal périodique présente un motif qui se répète identiquement après une durée minimale $T$, la période. La fréquence $f$ compte le nombre de motifs par seconde.
\begin{loi}$f=1/T$, avec $T$ en secondes et $f$ en hertz. Pour améliorer une lecture graphique, mesurer plusieurs périodes : si cinq périodes durent 10 ms, $T=2,0$ ms et $f=500$ Hz.\end{loi}
\schema[corrige]{signal-periodique}
Une demi-oscillation n'est pas une période. Choisir deux points dans le même état du motif, par exemple deux maxima successifs.
\essSection{2}{Hauteur et timbre}
À l'audition humaine, les sons audibles se situent approximativement entre 20 Hz et 20 kHz ; les limites dépendent des personnes. Au-dessous : infrasons ; au-dessus : ultrasons. Dans le domaine audible, une fréquence plus élevée correspond à un son plus aigu.
\schema[corrige]{timbres}
Deux sons de même fréquence peuvent avoir des formes de signaux différentes, donc des timbres différents. À forme et fréquence comparables, une amplitude plus importante est associée à une intensité sonore plus grande. La fréquence n'est pas une mesure de volume sonore.
\essSection{3}{Intensité et niveau}
Le niveau d'intensité sonore s'exprime en décibels. Son échelle n'est pas linéaire : on n'additionne pas les niveaux comme des longueurs. Dans les tableaux d'entraînement, une intensité multipliée par 10 correspond à une hausse de niveau de 10 dB ; on utilise cette correspondance fournie sans calcul logarithmique.
\begin{attention}[Exposition sonore]Le risque dépend du niveau et de la durée d'exposition, ainsi que de la proximité de la source. Réduire le niveau, s'éloigner et limiter la durée réduisent l'exposition. Une application de téléphone non étalonnée ne remplace pas une mesure de référence.\end{attention}
\essSection{4}{Expérimenter et modéliser}
\begin{experience}[Enregistrer et comparer]Enregistrer deux instruments jouant la même note à distance constante. Lire plusieurs périodes, comparer fréquence et forme, puis modifier doucement le niveau sans rechercher un son fort. Un microcontrôleur et un petit transducteur peuvent produire un signal de fréquence programmée ; identifier la chaîne commande, vibration, propagation, microphone et acquisition.\end{experience}
\exercices
\rubrique{Prendre confiance}
\begin{blocLivre}{ex-15-1}\begin{exercice}[1]{De la période à la fréquence}
Un signal a une période de 2,00 ms.
\begin{questions}
\item Exprimer $T$ en secondes.
\item Calculer $f$.
\end{questions}
\end{exercice}\end{blocLivre}
\begin{blocLivre}{ex-15-2}\begin{exercice}[1]{De la fréquence à la période}
Un diapason vibre à 250 Hz.
\begin{questions}
\item Calculer sa période.
\item Combien de motifs réalise-t-il en une seconde ?
\end{questions}
\end{exercice}\end{blocLivre}
\begin{blocLivre}{ex-15-3}\begin{exercice}[1]{Trois domaines}
On donne 12 Hz, 1000 Hz et 40 kHz.
\begin{questions}
\item Classer ces signaux dans les trois domaines.
\item Le signal de plus grande fréquence est-il nécessairement le plus intense ?
\end{questions}
\end{exercice}\end{blocLivre}
\begin{blocLivre}{ex-15-4}\begin{exercice}[1]{Lire un motif}
Utiliser le signal temporel fourni.
\schema{signal-periodique}
\begin{questions}
\item Lire la période.
\item En déduire la fréquence.
\end{questions}
\end{exercice}\end{blocLivre}
\rubrique{S’entraîner}
\begin{blocLivre}{ex-15-5}\begin{exercice}[2]{Mesurer plusieurs périodes}
Cinq périodes occupent 12,5 ms sur un enregistrement.
\begin{questions}
\item Calculer $T$ puis $f$.
\item Pourquoi cette méthode peut-elle améliorer la lecture ?
\end{questions}
\end{exercice}\end{blocLivre}
\begin{blocLivre}{ex-15-6}\begin{exercice}[2]{Deux instruments}
Les signaux A et B de la figure ont la même période.
\schema{timbres}
\begin{questions}
\item Comparer leurs hauteurs.
\item Comparer leurs timbres.
\end{questions}
\end{exercice}\end{blocLivre}
\begin{blocLivre}{ex-15-7}\begin{exercice}[2]{Plus aigu}
Deux sons audibles ont des fréquences de 220 Hz et 660 Hz.
\begin{questions}
\item Lequel est plus aigu ?
\item Calculer le rapport des périodes.
\end{questions}
\end{exercice}\end{blocLivre}
\begin{blocLivre}{ex-15-8}\begin{exercice}[2]{Une échelle en décibels}
Un tableau indique : intensité $I_0$ : 50 dB ; $10I_0$ : 60 dB ; $100I_0$ : 70 dB.
\begin{questions}
\item Que devient l’intensité entre 50 et 70 dB ?
\item Pourquoi 70 dB ne signifie-t-il pas une intensité 1,4 fois celle de 50 dB ?
\end{questions}
\end{exercice}\end{blocLivre}
\rubrique{Raisonner et approfondir}
\begin{blocLivre}{ex-15-9}\begin{exercice}[3]{Une confusion de paramètres}
Une application augmente l’amplitude d’un son enregistré sans changer son motif temporel.
\begin{questions}
\item La fréquence change-t-elle ?
\item Quelle sensation peut principalement changer ?
\end{questions}
\end{exercice}\end{blocLivre}
\begin{blocLivre}{ex-15-10}\begin{exercice}[3]{Un générateur programmé}
Un programme impose une période de 1,25 ms à un petit transducteur.
\begin{questions}
\item Déterminer la fréquence.
\item Proposer la période pour doubler cette fréquence.
\end{questions}
\end{exercice}\end{blocLivre}
\begin{blocLivre}{ex-15-11}\begin{exercice}[3]{Un son non périodique}
Un choc bref est enregistré ; aucun motif répétitif n’est identifiable.
\begin{questions}
\item Peut-on utiliser une période unique lue entre deux bosses quelconques ?
\item Pourquoi faut-il d’abord vérifier la répétition ?
\end{questions}
\end{exercice}\end{blocLivre}
\begin{blocLivre}{ex-15-12}\begin{exercice}[3]{Réduire une exposition}
Une personne se tient près d’une enceinte pendant un concert.
\begin{questions}
\item Citer trois moyens de réduire son exposition sonore.
\item Pourquoi le seul nom du style musical ne permet-il pas de juger le risque ?
\end{questions}
\end{exercice}\end{blocLivre}

